\usepackage[a4paper,
    mag=1000, includefoot,
    left=2.5cm, right=1.5cm, top=2cm, bottom=2cm,
    headsep=0.7cm, footskip=1cm
]{geometry}

\usepackage{mathtext}
\usepackage{amsmath,amssymb,amsthm,amscd,amsfonts}
\usepackage{cmap}
\usepackage[utf8]{inputenc}
\usepackage[T2A]{fontenc}
\usepackage[english,russian]{babel}
\usepackage{graphicx}
\usepackage[unicode,pdftex]{hyperref}
\usepackage{diagbox}
\usepackage{float}
\usepackage{tocloft}
\usepackage{subcaption}
\usepackage{algorithm}
\usepackage{algpseudocode}

\DeclareMathOperator*{\argmin}{arg\,min}
\DeclareMathOperator*{\argmax}{arg\,max}
\DeclareMathOperator{\rank}{rank}
\DeclareMathOperator{\NOKop}{НОК}
\newcommand{\ii}{\mathrm{i}}
\newcommand{\Hankel}{\mathcal H}

\newtheorem{statement}{Утверждение}[section]
\newtheorem{corollary}[statement]{Следствие}
\newtheorem{remark}[statement]{Замечание}

\algrenewcommand{\algorithmicrequire}{\textbf{Входные данные:}}
\algrenewcommand{\algorithmicensure}{\textbf{Выходные данные:}}

\graphicspath{
    {assets/images/chapter1/}
    {assets/images/chapter2/}
}

\hypersetup{
    pdftitle={Анализ сингулярного спектра для препроцессинга при обнаружении прямых линий на изображении},
    pdfauthor={Фалин Алексей Дмитриевич},
    colorlinks=true,
    linkcolor=black,
    citecolor=black,
    urlcolor=blue
}
